\documentclass[11pt,reqno]{amsart}
\usepackage[T1]{fontenc}
\usepackage{amsmath,amssymb,mathtools}
\usepackage[letterpaper,margin=1in]{geometry}
\usepackage{microtype}
\usepackage[hidelinks]{hyperref}
\hypersetup{pdftitle={Statement of Purpose: New Trends in Tropical Geometry},pdfauthor={Juliette Bruce}}
\setlength{\parindent}{0pt}
\setlength{\parskip}{5.5pt plus 1pt minus 1pt}
\setlength{\emergencystretch}{1em}
\pagestyle{plain}
\newcommand{\cA}{\mathcal{A}}
\newcommand{\cM}{\mathcal{M}}
\newcommand{\cX}{\mathcal{X}}
\newcommand{\Mbar}{\overline{\mathcal{M}}}
\newcommand{\Q}{\mathbb{Q}}
\newcommand{\C}{\mathbb{C}}
\newcommand{\Pp}{\mathbb{P}}
\newcommand{\Z}{\mathbb{Z}}
\newcommand{\trop}{\mathrm{trop}}
\DeclareMathOperator{\Gr}{Gr}
\DeclareMathOperator{\GL}{GL}

\newtheorem{theorem}{Theorem}

\begin{document}
\begin{center}
{\large\bfseries\scshape Juliette Bruce's Statement of Purpose}\\[3pt]
{\itshape New Trends in Tropical Geometry}
\end{center}
\vspace{-5pt}

My research uses combinatorial methods to study the topology of moduli spaces and questions of positivity in algebraic geometry. During the program \emph{New Trends in Tropical Geometry}, I plan to pursue two related directions: developing tropical methods to study moduli spaces of abelian varieties and using tropical geometry to study effective divisors, interpolation, and vector bundles. These projects build on my work in commutative algebra, toric geometry, and tropical geometry.

\textbf{Tropical Geometry and Moduli Spaces.}
A guiding principle in my work is that the combinatorics of degenerations can reveal geometry that is difficult to access directly. If $\cX$ is a smooth connected separated Deligne--Mumford stack over $\C$ with a normal crossings compactification $\cX\subset\overline{\cX}$, its boundary complex $\Delta$ computes the weight-zero compactly supported cohomology \cite[\S5]{CGP21}:
\[
W_0H_c^k(\cX;\Q)\cong\widetilde H^{k-1}(\Delta;\Q)\qquad(k>0).
\]
For $\cM_g$, this complex is the link of the tropical moduli space of curves. Chan, Galatius, and Payne used this perspective to produce large families of unstable cohomology classes. With Brandt, Chan, Melo, Moreland, and Wolfe, I developed the analogous approach for $\cA_g$, the moduli space of principally polarized abelian varieties \cite{BBCMMW24}. 

\begin{theorem}\cite[Thm.~A]{BBCMMW24}\label{thm:Ag}
The rational cohomology satisfies $H^{k}(\cA_{g};\Q)\neq0$ for $(g,k)=(5,15)$, $(5,20)$, $(6,30)$, $(7,28)$, $(7,33)$, $(7,37)$, and $(7,42)$.
\end{theorem}

Our computations produced new unstable cohomology of $\mathrm{Sp}_{2g}(\Z)$; in particular, $H^{15}(\cA_5;\Q)\neq0$ answers a question of Grushevsky about odd-degree cohomology. More recently, with Assaf, Brandt, Chan, and Vlad, I constructed tropicalizations of locally symmetric varieties whose links compute weight-zero compactly supported cohomology, producing new unstable cohomology classes for principal congruence subgroups of $\GL_g(\Z)$ \cite{ABBCV25}.

In separate work with the same collaborators, we introduced level structures on tropical abelian varieties and gave a modular interpretation of the resulting tropical moduli spaces \cite{ABBCV26}. Graph complexes and the matroidal geometry of $\cA_g$ also motivated my work with Bucciarelli and Zacovic on complexes of matroids \cite{BBZ26}. I plan to pursue the following projects.

\emph{Cone stacks.}
With Raluca Vlad, I am developing a cone-stack refinement of tropical moduli spaces of abelian varieties with level structure. Their present modular interpretation identifies points of a generalized cone complex with isomorphism classes; following the theory of tropical curves as cone stacks \cite{CCUW20}, we seek a moduli-theoretic formulation that retains automorphisms. This would allow level-changing maps and the tropical Torelli map to be studied as morphisms of moduli stacks and provide a foundation for tropical intersection theory on these spaces. I then hope to extend this construction to tropical locally symmetric varieties. A longer-term goal is a tropical proof of Mumford's proportionality theorem \cite{Mum77}.

\emph{Realizing weight zero.}
With Melody Chan and Dan Petersen, I am constructing a continuous proper map
$\lambda\colon\cA_g(\C)\to\cA_g^{\trop}$
whose pullback on compactly supported cohomology has image $W_0H_c^*(\cA_g;\Q)$. This is an analogue of Chan--Galatius--Payne's construction for curves using short hyperbolic geodesics \cite[\S7]{CGP21}. For abelian varieties, sufficiently short lattice vectors span an isotropic subspace for the polarization, providing the starting point for a tropical construction. Beyond this case, I hope to understand when the weight-zero comparison for a toroidal compactification can be realized by a proper map to its associated tropical cone complex.

\emph{Beyond weight zero.}
I also plan to study the weight-two compactly supported cohomology of $\cA_g$. Payne and Willwacher's work on $\cM_{g,n}$ suggests enriching the boundary complex with low-degree cohomology of boundary strata \cite{PW24}. For $\cA_g$, I aim to construct an analogous complex from toroidal boundary strata and their incidence maps, combining rational polyhedral cones with the cohomology of lower-dimensional moduli spaces. I hope to explore this with Amy Li, a new postdoc at Dartmouth. If successful, I hope this framework will allow me to extend the exponential growth in $g$ established in weight zero by Brown, Chan, Galatius, and Payne \cite{BCGP24} to higher weights.

I am also working to understand how Engel, de Gaay Fortman, and Schreieder use regular matroids and degenerations to produce counterexamples to the integral Hodge conjecture for abelian varieties \cite{EGFS25}. I hope to clarify how their obstruction to algebraicity relates to tropical geometry.

\textbf{Tropical Geometry and Positivity.}
A second direction of my research uses tropical compactifications and degenerations to study divisors and cohomology. With Cavalieri, Maclagan, and Huang, I am studying the effective cone of $\Mbar_{0,n}$ and related notions of positivity. The realization of $\Mbar_{0,n}$ as a tropical compactification gives access to rich combinatorial structures \cite{GM10}. We aim to use tropical cycles, boundary strata, and intersection theory to obtain new effective classes and constraints on positivity. I hope to explain tropically why the pseudoeffective cone is polyhedral for $n\leq6$ \cite{HT02,Cas09} and nonpolyhedral for $n\geq8$ \cite{CLTU23,Mul25}, and to use this perspective to approach the remaining case $n=7$.

Kapranov's realization of $\Mbar_{0,n}$ as an iterated blow-up of $\Pp^{n-3}$ \cite{Kap93} connects these questions about effective divisors to interpolation on projective space and, more broadly, on toric varieties. I plan to use tropical degenerations of point configurations to study linear systems of plane curves with prescribed multiplicities. The aim is to obtain combinatorial bounds on their dimensions and identify configurations for which these bounds are sharp. Nagata-type inequalities and the more refined predictions of the Segre--Harbourne--Gimigliano--Hirschowitz conjecture \cite{CHMR13} provide concrete tests for these methods.

With Chris Manon, I am developing a theory of higher cohomology for tropical toric vector bundles. Building on Klyachko's classification \cite{Kly90} and the tropical frameworks of Khan--Maclagan \cite{KhM26} and Kaveh--Manon \cite{KM24}, we aim to construct suitable notions of matroid-valued sheaves on the fan of a toric variety. The goal is to define higher cohomology for these tropical toric vector bundles, including those with nonrealizable underlying matroids, with familiar nice properties of cohomology of vector bundles on varieties. I am particularly interested in cohomological vanishing criteria and their consequences for multigraded Castelnuovo--Mumford regularity \cite{MS04}, connecting with my work on homological algebra over toric varieties. More speculatively, this  would serve as a first small step in the direction of my dream of a more richly developed connections between tropical geometry and commutative algebra. 

\textbf{Relevance of the Visit and Professional Goals.}
The program's emphasis on tropical moduli spaces and matroids aligns closely with these projects. Sustained discussions with researchers working on tropical intersection theory, moduli, and n son-Archimedean geometry would help me advance all of the ongoing work described above. In particular, I would like to understand how the proposed realization map relates to non-Archimedean tropicalization and Berkovich skeleta \cite{Viv13,ACP15}.

My work has long been rooted in combinatorial algebraic geometry, and I began working deeply in tropical geometry as a postdoc. A semester at SLMath would allow me to become more fully part of this research community, strengthen existing collaborations, and develop new ones. Exposure to the program's broader themes, including real tropical geometry and tropical mirror symmetry, would help me formulate questions beyond the parts of the subject I currently know best.

I would also contribute to the program through working seminars and mentoring graduate students and postdocs. My experience organizing the GEMS workshops and serving as the inaugural president of Spectra has shaped my commitment to supporting mathematicians and building welcoming mathematical communities. I would bring that same commitment to SLMath. I would be enthusiastic about creating a vibrant and engaging social and working environment throughout the program. 

\clearpage
\begingroup
\small
\setlength{\parskip}{0pt}
\begin{thebibliography}{ABBCV25}
\setlength{\itemsep}{5pt}

\bibitem[ACP15]{ACP15}
D.~Abramovich, L.~Caporaso, and S.~Payne,
\emph{\href{https://www.numdam.org/item/ASENS_2015__48_4_765_0/}{The tropicalization of the moduli space of curves}},
Ann. Sci. \'Ec. Norm. Sup\'er. (4) \textbf{48} (2015), no.~4, 765--809.

\bibitem[ABBCV25]{ABBCV25}
E.~Assaf, M.~Brandt, J.~Bruce, M.~Chan, and R.~Vlad,
\emph{\href{https://arxiv.org/abs/2505.10504}{Tropicalizations of locally symmetric varieties}},
arXiv:2505.10504 (2025).

\bibitem[ABBCV26]{ABBCV26}
E.~Assaf, M.~Brandt, J.~Bruce, M.~Chan, and R.~Vlad,
\emph{\href{https://arxiv.org/abs/2609.14895}{Level structures on tropical abelian varieties}},
arXiv:2609.14895 (2026).

\bibitem[BBC$^+$24]{BBCMMW24}
M.~Brandt, J.~Bruce, M.~Chan, M.~Melo, G.~Moreland, and C.~Wolfe,
\emph{\href{https://doi.org/10.2140/gt.2024.28.497}{On the top-weight rational cohomology of $\cA_g$}},
Geom. Topol. \textbf{28} (2024), no.~2, 497--538.

\bibitem[BCGP24]{BCGP24}
F.~Brown, M.~Chan, S.~Galatius, and S.~Payne,
\emph{\href{https://arxiv.org/abs/2405.11528}{Hopf algebras in the cohomology of $\cA_g$, $\mathrm{GL}_n(\Z)$, and $\mathrm{SL}_n(\Z)$}},
arXiv:2405.11528 (2024).

\bibitem[BBZ26]{BBZ26}
J.~Bruce, J.~Bucciarelli, and B.~Zacovic,
\emph{\href{https://arxiv.org/abs/2605.24695}{Explorations of matroid complexes}},
arXiv:2605.24695 (2026).

\bibitem[Cas09]{Cas09}
A.-M.~Castravet,
\emph{\href{https://doi.org/10.1090/S0002-9947-09-04641-8}{The Cox ring of $\Mbar_{0,6}$}},
Trans. Amer. Math. Soc. \textbf{361} (2009), no.~7, 3851--3878.

\bibitem[CLTU23]{CLTU23}
A.-M.~Castravet, A.~Laface, J.~Tevelev, and L.~Ugaglia,
\emph{\href{https://doi.org/10.1515/crelle-2023-0022}{Blown-up toric surfaces with non-polyhedral effective cone}},
J. Reine Angew. Math. \textbf{800} (2023), 1--44.

\bibitem[CCUW20]{CCUW20}
R.~Cavalieri, M.~Chan, M.~Ulirsch, and J.~Wise,
\emph{\href{https://doi.org/10.1017/fms.2020.16}{A moduli stack of tropical curves}},
Forum Math. Sigma \textbf{8} (2020), Paper No.~e23.

\bibitem[CGP21]{CGP21}
M.~Chan, S.~Galatius, and S.~Payne,
\emph{\href{https://doi.org/10.1090/jams/965}{Tropical curves, graph complexes, and top weight cohomology of $\cM_g$}},
J. Amer. Math. Soc. \textbf{34} (2021), no.~2, 565--594.

\bibitem[CHMR13]{CHMR13}
C.~Ciliberto, B.~Harbourne, R.~Miranda, and J.~Ro\'e,
\emph{\href{https://arxiv.org/abs/1202.0475}{Variations on Nagata's conjecture}},
Clay Math. Proc. \textbf{18} (2013), 185--203.

\bibitem[EGFS25]{EGFS25}
P.~Engel, O.~de Gaay Fortman, and S.~Schreieder,
\emph{\href{https://arxiv.org/abs/2507.15704}{Matroids and the integral Hodge conjecture for abelian varieties}},
arXiv:2507.15704 (2025).

\bibitem[GM10]{GM10}
A.~Gibney and D.~Maclagan,
\emph{\href{https://doi.org/10.2140/ant.2010.4.855}{Equations for Chow and Hilbert quotients}},
Algebra Number Theory \textbf{4} (2010), no.~7, 855--885.

\bibitem[HT02]{HT02}
B.~Hassett and Y.~Tschinkel,
\emph{\href{https://arxiv.org/abs/math/0110231}{On the effective cone of the moduli space of pointed rational curves}},
Contemp. Math. \textbf{314} (2002), 83--96.

\bibitem[Kap93]{Kap93}
M.~M.~Kapranov,
\emph{\href{https://doi.org/10.1090/advsov/016.2/02}{Chow quotients of Grassmannians.~I}},
in \emph{I.~M.~Gelfand Seminar},
Adv. Soviet Math. \textbf{16}, Part~2, Amer. Math. Soc., Providence, RI, 1993, 29--110.

\bibitem[KM24]{KM24}
K.~Kaveh and C.~Manon,
\emph{\href{https://arxiv.org/abs/2405.03576}{Tropical vector bundles and matroids}},
arXiv:2405.03576 (2024).

\bibitem[KhM26]{KhM26}
B.~Khan and D.~Maclagan,
\emph{\href{https://doi.org/10.1016/j.aim.2025.110646}{Tropical toric vector bundles}},
Adv. Math. \textbf{484} (2026), Paper No.~110646.

\bibitem[Kly90]{Kly90}
A.~A.~Klyachko,
\emph{\href{https://doi.org/10.1070/IM1990v035n02ABEH000707}{Equivariant bundles on toral varieties}},
Math. USSR-Izv. \textbf{35} (1990), no.~2, 337--375.

\bibitem[MS04]{MS04}
D.~Maclagan and G.~G.~Smith,
\emph{\href{https://doi.org/10.1515/crll.2004.040}{Multigraded Castelnuovo--Mumford regularity}},
J. Reine Angew. Math. \textbf{571} (2004), 179--212.

\bibitem[Mul25]{Mul25}
S.~Mullane,
\emph{\href{https://arxiv.org/abs/2510.25044}{Isoresidual fibrations and the birational geometry of moduli of pointed stable curves}},
arXiv:2510.25044 (2025).

\bibitem[Mum77]{Mum77}
D.~Mumford,
\emph{\href{https://www.dam.brown.edu/people/mumford/alg_geom/papers/1977b--HirzProp-GT.pdf}{Hirzebruch's proportionality theorem in the non-compact case}},
Invent. Math. \textbf{42} (1977), 239--272.

\bibitem[PW24]{PW24}
S.~Payne and T.~Willwacher,
\emph{\href{https://doi.org/10.1215/00127094-2024-0003}{Weight two compactly supported cohomology of moduli spaces of curves}},
Duke Math. J. \textbf{173} (2024), no.~16, 3107--3178.

\bibitem[Viv13]{Viv13}
F.~Viviani,
\emph{\href{https://arxiv.org/abs/1204.3875}{Tropicalizing vs compactifying the Torelli morphism}},
in \emph{Tropical and Non-Archimedean Geometry},
Contemp. Math. \textbf{605} (2013), 181--210.

\end{thebibliography}
\endgroup
\end{document}
